\documentclass[
pdfa=false,  % to avoid compilation error
color=9c,
logo=body,
class=article,
marginpar=false,
logofile=../logos/tuda_logo.pdf
]{tudapub}

% ---------- Packages ---------- % 

\usepackage[main=english, ngerman]{babel}
\usepackage{lipsum} % for dummy text
\usepackage{subfiles} % Best loaded last in the preamble
\usepackage{caption} % To center captions
\usepackage{tikz} % For drawing
\usepackage{float} % To fix figure positions where they should be...
\usetikzlibrary{external} % To externalize tikz figures
\usetikzlibrary{arrows} % For fancy errors in tikz figures
\usepackage{amsmath} % Use text in math mode
\usepackage[toc,page]{appendix} % For appendix

% ---------- Notes ---------- % 

% TODO: Do we want two columns?


% ---------- Make Title ---------- % 

\author{Yannik P. Frisch, Maximilian A. Gehrke}
\title{Predicting Task Performance - How to Quantify Metacognition Using Drawn Confidence Distributions}
\date{March 10, 2020}
\subtitle{Applied Cognitive Science}
\addTitleBoxLogo*{\includegraphics[width=.7\linewidth]{../logos/ccs_logo.png}}


\begin{document}
\maketitle

% ---------- Abstract ---------- % 
\begin{abstract}
	Several studies suggest that humans are rather poor when judging their own performance after fulfilling a task. We present a new method using drawn probability density functions to measure subjects' confidence of their performance on various tasks of a questionnaire that we designed specifically for this purpose. We then evaluate their responses using an evaluation function unknown to the participants and assess their confidence by statistically measuring the discrepancy between their expected and the actual task performance. We conclude that people are poor at assessing their own task performance, overestimate poor task performance and underestimate good task performance, and show a cognitive anchoring bias towards an average performance rating. We suggest further experiments on this topic, using a greater variety of tasks and a higher number of subjects.
	
\end{abstract}

\section{Introduction (Yannik P. Frisch)}
	\label{sec:introduction}
	\subfile{../sections/introduction.tex}

\newpage
\section{Method (Maximilian A. Gehrke)}
	\label{sec:method}
	\subfile{../sections/method.tex}

\newpage
\section{Results (Yannik P. Frisch)}
	\label{sec:results}
	\subfile{../sections/results.tex}

\newpage
\section{Discussion (Maximilian A. Gehrke)}
	\label{sec:discussion}
	\subfile{../sections/discussion.tex}


% ---------- Bibliography ---------- %
\newpage
\bibliography{../bib/references.bib}{}
\bibliographystyle{plain}


% ---------- Appendix ---------- %
\newpage
\begin{appendices}
	\section{The Questionnaire}
		\label{appendix:questionnaire}
		\subfile{../sections/questionnaire.tex}
	\newpage
	\section{Individual Results}
		\label{appendix:individual_results}
		\subfile{../sections/individual_results.tex}
\end{appendices}


\end{document}